% 正文样例：参考其表达与结构，内容和星级以实际课程来源为准。
\section{平台与网络效应}

\StudyTerm{\PriorityImportant}{平台}{组织参与者互动、提供互动条件并制定和执行规则的中介。}
\textbf{四类交易安排：}收费改变参与和交易成本；搜索、推荐或派单改善匹配；评价降低信息不对称；取消、退款和纠纷规则分配履约风险。会员费对应接入或一定时期的使用，交易费随成交收取。

\StudyTerm{\PriorityMust}{网络效应}{新增用户改变其他用户从产品、服务或互动中获得的价值，属于消费外部性。}

\StudyTerm{\PriorityMust}{组内网络效应}{同一群体的用户相互影响，如微信联系人增加提高其他用户的通信价值。}

\StudyTerm{\PriorityMust}{跨组网络效应}{一类用户的规模或活动改变另一类用户的价值，如买家增加吸引卖家、卖家增加丰富买家选择。}

\StudyTerm{\PriorityImportant}{正向与负向网络效应}{按对其他用户价值的影响方向区分。互补可带来正向效应，拥挤、干扰或资源竞争可带来负向效应。}
\textbf{两个分类维度独立：}骑手争夺有限订单属于负向组内效应，消费者增加吸引商家属于正向跨组效应。

\StudyTerm{\PriorityImportant}{规模经济}{产量增加使平均资源成本下降，作用于供给侧；网络效应改变用户价值，作用于需求侧。}

\StudyPoint{\PriorityImportant}{场景分化与多个稳定大网络}\label{kp:scene-networks}
\textbf{前提：}不同场景的网络效应范围相对独立，可近似分开分析。
$n_s$ 为启动临界规模，$n_l$ 为稳定大网络规模。
% 原截图的连续文字推理明显超过100字；短段无需仿照画图。
{\setlength{\topsep}{2pt}
\begin{center}
\begin{tikzpicture}[study logic base]
  \node (scope) {使用\textbf{场景分化}\\效应范围相对独立};
  \node[right=of scope] (local) {各自形成\textbf{预期}\\分别确定 $n_s$ 与 $n_l$};
  \node[right=of local,text width=4.1cm] (large) {若各自形成\textbf{稳定大网络}\\则多个大网络\textbf{可并存}};
  \draw[study logic arrow] (scope) -- (local);
  \draw[study logic arrow] (local) -- (large);
\end{tikzpicture}
\end{center}
}
\begin{center}
{\small
\begin{tabular}{ll}
\toprule
资料中的场景示例 & 网络效应范围的示意口径\\
\midrule
微信支付 & 社交关系、周边商户及线下支付\\
支付宝 & 电商交易、金融与公共服务入口\\
\bottomrule
\end{tabular}}
\end{center}
\begin{itemize}
\item \textbf{范围内多重均衡：}给定效应范围及满足条件的接入费，可能有多个均衡；预期影响实际实现的规模。
\item \textbf{边界：}跨范围的大网络并存不等于范围内多重均衡；仅凭范围重合也不能推出胜者通吃。
\end{itemize}

\section{双边收费}

\StudyTerm{\PriorityMust}{价格结构非中性}{两侧总费用不变时，改变买卖双方的收费分配仍会改变交易量。}
按成交向两侧分别收取 $p^B,p^S$；潜在配对总量归一化，两侧收益独立。
$D^B,D^S$ 表示各侧愿意交易的比例，均随本侧收费提高而下降，故成交量为 $Q=D^B(p^B)D^S(p^S)$。此处假定两侧不能通过转移支付抵消收费分配。

\StudyPoint{\PriorityMust}{固定总价下的最优收费结构}
\begin{proofbox}[title=推导]
固定总费用 $p=p^B+p^S$，两侧需求可微且为正。设 $q(p^B)=D^B(p^B)D^S(p-p^B)$。
由于 $\mathrm dp^S/\mathrm dp^B=-1$，
\begin{align*}
q'&=(D^B)'D^S-D^B(D^S)',\\
\frac{q'}q&=\frac{(D^B)'}{D^B}-\frac{(D^S)'}{D^S}
=h_S-h_B,\qquad h_i=-\frac{(D^i)'}{D^i}>0.
\end{align*}
$h_i$ 为本侧收费增加一单位引起的需求相对降幅。
\textbf{内点最优的必要条件：}$h_B=h_S$，仍须检验最优性并比较可行边界。
若 $h_B>h_S$，小幅降低买方收费、等额提高卖方收费可增加交易量。
\end{proofbox}
\begin{proofbox}[title=局部平稳与全局中性]
设 $0<p<2$，以下比较在线性需求均为正的可行区间内进行。
\begin{align*}
D^B=1-p^B,\quad D^S=1-p^S
&\Rightarrow q=(1-p^B)(1-p+p^B),\\
q'=p-2p^B,\quad q''=-2
&\Rightarrow q'(p/2)=0\quad\text{但 }q\text{ 不恒定};\\
D^i=e^{-ap^i},\quad a>0
&\Rightarrow Q=e^{-a(p^B+p^S)}=e^{-ap}\quad\text{只依赖总价}.
\end{align*}
\textbf{区别：}某点一阶变化为零 $\neq$ 所有可行收费分配下交易量不变。
\end{proofbox}

% 对应例题就近归入同一模型；名称与星级放框内题头，省去重复入口。
\begin{examplebox}[title={\PriorityMust\quad 例题：非负收费约束}]
\textbf{题：}总价 $p=0.4$，$D^B=1-2p^B$、$D^S=1-p^S$，两侧收费非负。求交易量最大的收费结构。

\textbf{解：}$p^B\in[0,0.4]$，$p^S=0.4-p^B$。
\begin{align*}
q&=(1-2p^B)(0.6+p^B)=0.6-0.2p^B-2(p^B)^2,\\
q'&=-0.2-4p^B,\qquad q''=-4.
\end{align*}
驻点 $p^B=-0.05$ 不可行；在可行区间 $q'<0$，故左端点最优：
\[
\boxed{p^{B*}=0,\qquad p^{S*}=0.4,\qquad Q^*=0.6.}
\]
\end{examplebox}

\section{新旧网络竞争：需求与收费}

\StudyTerm{\PriorityMust}{满足预期均衡（FEE）}{用户依据对网络规模的预期作出参与决策，由这些决策形成的实际规模恰好等于预期规模。若参与映射为 $n=g(n^e)$，则均衡同时满足 $n=n^e$，即 $n=g(n)$。它要求预期与行为结果一致，不意味着均衡唯一或稳定；同一价格可能对应多个满足预期均衡，稳定性需另按参与调整过程判断。}\label{kp:fee}

\StudyPoint{\PriorityMust}{新平台的需求与内点最优收费}\label{kp:new-network}
\textbf{设定：}消费者均匀分布于 $x\in[0,1]$，单归属，市场完全覆盖。新平台在左端、旧网络在右端；旧网络免费，边际成本为零。
\begin{center}
{\small\begin{tabular}{ll}
\toprule
符号 & 含义与条件\\
\midrule
$r$ & 共同基础收益\\
$q$ & 新平台质量优势\\
$\tau>0$ & 单位偏好距离成本\\
$\beta_i\ge0$ & 各网络的组内效应强度\\
$A_1\ge0$ & 新平台接入费\\
\bottomrule
\end{tabular}}
\end{center}
% 课件中的线性城市分界关系：保留位置、边际用户和份额，不扩写图注。
\begin{center}
\begin{tikzpicture}[x=9cm,every node/.style={font=\small}]
  \draw[thick] (0,0) -- (1,0);
  \fill (0,0) circle (1.3pt) node[above] {新平台 1};
  \fill (1,0) circle (1.3pt) node[above] {旧网络 2};
  \fill[starred] (0.55,0) circle (1.5pt) node[above] {$\hat x$};
  \node[below] at (0,0) {$0$};
  \node[below] at (1,0) {$1$};
  \node[below] at (0.275,-0.2) {$n_1=\hat x$};
  \node[below] at (0.775,-0.2) {$n_2=1-\hat x$};
  \node[below,font=\footnotesize] at (0.5,-0.65) {分界关系示意，位置不代表具体数值};
\end{tikzpicture}
\end{center}
\begin{proofbox}[title=需求与利润的推导]
两种选择的效用为
\begin{align*}
u_1(x)&=r+q+\beta_1n_1^e-\tau x-A_1,\\
u_2(x)&=r+\beta_2n_2^e-\tau(1-x).
\end{align*}
无差异用户满足 $u_1(\hat x)=u_2(\hat x)$。由于消费者均匀分布且 $\tau>0$，$n_1=\hat x,n_2=1-\hat x$。按\hyperref[kp:fee]{FEE}取 $n_i^e=n_i$：
\begin{align*}
q+\beta_1\hat x-\tau\hat x-A_1
  &=\beta_2(1-\hat x)-\tau(1-\hat x),\\
D\hat x&=q+\tau-\beta_2-A_1,\qquad D=2\tau-\beta_1-\beta_2>0,\\
n_1(A_1)&=\frac{C-A_1}{D},\qquad C=q+\tau-\beta_2.
\end{align*}
$D>0$ 保证该分割支路需求向下；还须 $0<n_1<1$。利润与最优条件为
\begin{align*}
\Pi_1(A_1)&=\frac{A_1(C-A_1)}D,\\
\Pi_1'(A_1)&=\frac{C-2A_1}D=0,\qquad
\Pi_1''(A_1)=-\frac2D<0,\\
A_1^*&=\frac C2,\qquad n_1^*=\frac C{2D},\qquad
\Pi_1^*=\frac{C^2}{4D}.
\end{align*}
\textbf{内点有效范围：}$0<C<2D$，即
$\beta_2-\tau<q<3\tau-2\beta_1-\beta_2$。
区间外不能直接套内点公式，应改用份额在 $[0,1]$ 内的需求比较边界。
\end{proofbox}
\begin{examplebox}[title={\PriorityImportant\quad 例题：新平台收费}]
\textbf{题：}$\tau=1$、$\beta_1=\beta_2=0.6$、$q=0.5$，其余条件见\hyperref[kp:new-network]{本模型设定}，求新平台最优收费、份额与利润。

\textbf{解：}$D=2-0.6-0.6=0.8>0$，$C=0.5+1-0.6=0.9$，
且 $0<0.9<1.6=2D$，满足内点条件。
\[
A_1^*=\frac{0.9}{2}=0.45,\qquad
n_1^*=\frac{0.9}{1.6}=0.5625,\qquad
\Pi_1^*=\frac{0.81}{3.2}=0.253125.
\]
\end{examplebox}
